
\documentclass[11pt]{article}

\usepackage[margin=1in]{geometry}
\usepackage{amsmath,amssymb}
\usepackage[T1]{fontenc}
\usepackage{enumitem}
\usepackage{silence}
\usepackage{apacite}

\title{Review: A Unified Blister and Subglacial Hydrology Framework for Supraglacial Lake Drainage Events}
\author{}
\date{}

\begin{document}

\maketitle

\section*{Summary}

The authors have further developed an existing subglacial hydrology model to include blister
dynamics in addition to meltwater flow through the distributed and channelized components of
the modeling framework. Subglacial blisters form following supraglacial lake drainage events,
which are observed across Greenland. The rapid input of water to the bed causes localized uplift
of the glacier and can alter the ice flow velocity, both of which are observed in satellite and field
observations.

In the new modeling framework, blister dynamics are controlled by effective viscosity and a
leakage length scale. The introduction of these terms to the PDE system describing water flow
through the subglacial environment allows for water from the blister to enter into the distributed
and channelized components of the system, and also for the blister itself to propagate down
glacier introducing a new meltwater pathway in the model. The parameter space of the new
modeling framework is explored using an idealized setup applied to both a winter and
summertime representation of the subglacial environment. Additionally, the model is applied to a
regional case study, and is able to produce dynamics consistent with observations following a
lake drainage event in western Greenland.

\section*{Overall Statement}

Generally speaking, this manuscript is well constructed, well written, and the work described
represents a significant development of subglacial hydrology modeling. I believe that it warrants
publication and my suggested revisions are all quite minor.

\section*{General Comments}

The introduction section could use more description of the physics relevant to this work and
more specific/detailed definitions. As it is currently written, the introduction is quite short and to
the point. As it is also well written, this is not a bad thing, but I believe that it would benefit the
reader to have more context regarding the physics governing water flow through the subglacial
environment and how this connects to ice flow dynamics. I would also suggest spending some
time here providing more detail about the field and satellite observations that provide evidence of
blister dynamics. These additions would help to emphasize why the modeling development here
by the authors, which is quite a substantial advancement, is addressing a significant knowledge
gap.

The methods section does not clearly indicate which equations in the model description are
established in previous published work on blister modeling versus novel components of this
subglacial hydrology modeling framework with incorporated blister physics. It is stated that the
new terms linking the existing subglacial hydrology model to the blister dynamics are $\Sigma_m$
and $Q_m$. However, there are many more equations in this section, and it’s not entirely clear if they all
are based on Hewitt (2013), are novel to this work, or are pulling from additional sources. The
work presented here is very interesting and is a useful advancement for the capabilities of
subglacial hydrology models to better represent observations. However, this could be more
clearly emphasized by highlighting the new components here and how they connect to
established modeling frameworks. Additional citations throughout this section could also be
helpful in this distinction.

\section*{Specific Comments}

\subsection*{Key Points}

Currently stated key points are clear and accurate, but they do not emphasize the main
achievements of the work. I suggest reworking to highlight the advancement of subglacial
modeling capabilities represented in this work, and the new modeling framework's ability to
reproduce field observations. Shift the focus from the methods to the results. I would suggest
specifically mentioning the ability of the model to reproduce observed behavior of changes of
glacier dynamics following lake drainage events.

\subsection*{Abstract}

Line 26: Suggested addition to highlight differences: propagate $\rightarrow$ propagate slowly

\subsection*{Plain Language Summary}

Very clear and nicely written.

\subsection*{Introduction}

\begin{itemize}[leftmargin=*]
\item Line 54: Perhaps it is worth mentioning that supraglacial lake drainage events have been
observed in both summer and winter seasons (e.g. Dean et al., 2026). If the reader does not
already know this, it is not clear why there would be observed blister propagations in winter.

\item Line 56: What type of length scale do blisters propagate over: 100s of meters, kilometers?

\item Line 63: Generally a more detailed explanation of the physics of blisters in the subglacial system
would be useful in this first paragraph. The introduction quickly jumps from a quick overview of
why blisters are important to previous modeling efforts of blister dynamics. I think that the rest
of the paper would be well served if a bit more time was spent in the introduction highlighting
the gap that this work is filling in our ability to model the subglacial environment. Perhaps a
summary of field/satellite observations providing evidence of blister dynamics and/or a more
fleshed out explanation of the dynamics between ice flow and the subglacial environment could
help to highlight why it is important to incorporate blisters into subglacial hydrology models.

\item Line 76: I suggest adding a final sentence to this paragraph noting that these models do not
include a sophisticated representation of other components of the subglacial water flow, or note
if any do.

\item Line 83: This sentence gets to the heart of the paper. I think the point could be made stronger if a
bit more explanation was done early on in the introduction regarding why we should care about
blister dynamics and therefore their incorporation into subglacial hydrology models.
\end{itemize}

\subsection*{Method}

\begin{itemize}[leftmargin=*]
\item Line 165: Do the necessary observations to constrain this value exist or is the $h_0$ parameter
currently unknown and will require new/more field observations to determine?

\item Line 172: ``the Poisson’s ratio'' is a bit awkward, suggest changing to ``Poisson’s ratio'' or ``the
Poisson ratio''
\end{itemize}

\subsection*{Results}

\begin{itemize}[leftmargin=*]
\item Line 220--223: This is a bit confusing. Are channels allowed to grow in this model setup and they
have not developed because the model is representing a wintertime subglacial drainage system,
or are they intentionally prevented from evolving? After saturating the local cavity sheet, the
blister could initiate channel growth, but it doesn’t? I would think that 10 days is enough time to
see some initiation of channel growth, even if most of the melt remains within the blister.

\item Line 238: This sentence is a bit awkward, and it is long. I suggest breaking it in two.

\item Line 251: Super cool!

\item Line 253: Do the number of channels in the network influence the blister dynamics? No need to run
any additional tests, I would just be curious to know more if this is something the authors already
looked into.
\end{itemize}

\subsection*{Discussion}

\begin{itemize}[leftmargin=*]
\item Line 272: It would be helpful to remind the reader here what $\kappa$ represents in the lengthscale
equation.

\item Line 291: Effective viscosity is mentioned in the abstract, and then not again until here except in
equation (3) of the methods. Since it is an important parameter in governing blister propagation,
it would be helpful to go into more detail describing what the parameter is and how its value is
determined, either here or in the methods when it is first introduced.

\item Line 298: Is this range pulled from the literature?

\item Line 304: What is a bending dominated regime?

\item Line 310: Is this inferred value range consistent with previous parameterizations of effective
viscosity?

\item Line 318: It would be helpful to define bending-dominated versus gravity-dominated regimes.

\item Line 311--318: This paragraph could use more detail in general. The results are interesting, but
this explanation feels rushed.

\item Line 325: I had assumed that leakage referenced water moving from the blister to the distributed
component of the subglacial drainage system, but it seems that it also includes movement into
the channelized system. A definition would make things more clear.
\end{itemize}

\subsection*{Regional Study}

Line 390: Cool!

\subsection*{Conclusion}

Concise and nicely written.

\subsection*{Figures}

\begin{itemize}[leftmargin=*]
\item Figure 1: It would be helpful to reference this figure in the introduction text when explaining
what blisters are.

\item Figure 2: Really cool figure! Panels (b--d) Water flux magnitude is listed as the spatially mapped
variable in the background of the panels. It would be helpful to explicitly state which water flux
this is as several are referenced in the text and plotted previous panel. Total water flux in the
subglacial environment? Panel (f) displays the blister thickness at different locations over time.
This is very helpful to have, but would be more clear with a more detailed explanation of the
setup. I think a brief explanation in the text (line 225) of what the hypothetical monitoring
stations are would be useful.

\item Figure 3: Final sentence in caption is a bit awkward. Suggest removing.

\item Figure 6: Very nice figure!
\end{itemize}

\subsection*{Appendix A}

No comments.

\subsection*{Appendix B}

Line 496: Are these cases chosen because they represent end members of possible behavior of
the system?

\subsection*{Appendix C}

No comments.

\section*{References}

\begin{itemize}
\item Dean, C.W., Scharien, R., Willis, I. and McDougall, K.A., 2025. \textit{A decade of winter supraglacial
lake drainage across Northeast Greenland using C-band SAR}. EGUsphere, 2025, pp.1--38.

\item Hewitt, I.J., 2013. \textit{Seasonal changes in ice sheet motion due to melt water lubrication}.
Earth and Planetary Science Letters, 371, pp.16--25.
\end{itemize}

\end{document}
